\chapter{The Origin of Caste}

\section{Caste, Varna \& the Theories of Origin}

The basic unit of Paper 2 is caste. To understand the meaning of caste we must understand its origin; to understand the origin we must study the theories of origin.

\begin{CriticalPoint}
\textbf{None of the theories of the origin of caste is empirically valid} --- they are all based on speculation (``hava bazi''). There is no ``good'' theory here; every theorist is, in the teacher's words, ``a bad person.''
\end{CriticalPoint}

Six theories are dealt with: (1) the Traditional/Divine theory; (2) Karma and transmigration of soul; (3) the Racial theory; (4) the Occupational theory; (5) the Political theory; (6) the Broken Man theory (Ambedkar).

After origin, the next challenge is \textbf{Varna}. Caste, Varna and Jati must be separated (Senart; M. N. Srinivas). Notably, every thinker except M. N. Srinivas began with the idea of the origin of caste.

\begin{QuickRevision}
\begin{itemize}
\item Six theories of origin; none is empirically valid (all speculation).
\item Chain: caste, its origin, and the theories of origin.
\item Every thinker except M. N. Srinivas begins with the origin of caste.
\end{itemize}
\end{QuickRevision}

\section{The Traditional (Divine) Theory}

A creature called \textbf{Brahma} (the creator of the universe; the equivalent of the ``Big Bang'' in natural science) collapsed and fragmented into pieces. The four Varnas emerged from his body:
\begin{itemize}
\item From the \textbf{mouth} came the \textbf{Brahmins} --- the head/mouth symbolises knowledge, wisdom and good speech.
\item From the \textbf{arms} came the \textbf{Kshatriyas} --- arms symbolise strength and the warrior class.
\item From the \textbf{thighs} came the \textbf{Vaishyas} --- no scientific reason is given; merchants move about a great deal.
\item From the \textbf{feet} came the \textbf{Shudras} --- as the feet bear the body's burden, so the Shudras bear society's burden as labourers.
\end{itemize}

\begin{KeyConcept}
Conclusion of the divine theory: \textbf{caste has its origin in divinity} --- the deity made it, and we can do nothing about it.
\end{KeyConcept}

\begin{QuickRevision}
\begin{itemize}
\item Brahma's body: mouth: Brahmins; arms: Kshatriyas; thighs: Vaishyas; feet: Shudras.
\item Conclusion: caste has a divine origin.
\end{itemize}
\end{QuickRevision}

\section{Karma \& Transmigration of Soul}

Why is one born a Kshatriya? In this theory it is not Brahma but \textbf{karma} --- the good deeds of a past life are rewarded in the present. The transmigration of the soul means the fruit of past deeds is received in this life.

\begin{KeyConcept}
The \textbf{upper status} of a Varna is justified by \textbf{good karma} in a past life; the \textbf{lower status} is justified by \textbf{bad karma} (sin) in a past life.
\end{KeyConcept}

\begin{QuickRevision}
\begin{itemize}
\item Varna status is explained by karma in a past life (not by Brahma).
\item Good karma gives the upper Varna; bad karma gives the lower Varna.
\end{itemize}
\end{QuickRevision}

\section{Broken Man Theory (Ambedkar)}

This theory is given by \textbf{Ambedkar} (associated with his book ``The Broken Men''). The \textbf{Dalit} literally means the ``broken man'' (toota hua admi) --- broken out of the Varna vyavastha.

The argument runs:
\begin{enumerate}
\item The \textbf{original kings were Shudras} (the real kings), not Kshatriyas.
\item These Shudra kings \textbf{exploited the Brahmanical priests}.
\item The Brahmins held one power --- to \textbf{invest someone with the sacred thread}.
\item Enraged, the Brahmins refused the Shudras the sacred thread, so the Shudras became the only group without it and thereby got the \textbf{lowest status} among the four.
\end{enumerate}

\begin{ExampleBox}
\textbf{Jyotiba Phule} held the same opinion --- kings like \textbf{Bali} and \textbf{Sugriva} were Shudras.
\end{ExampleBox}

\subsection{The sacred thread and the Dvija--Ekaja divide}

The sacred thread ceremony (in Sanskrit, the \textbf{Upanayana} ceremony) is performed at about 7--8 years of age and marks a \textbf{spiritual rebirth} (divya janma). The thread itself is the ``dhaga.''

\begin{KeyConcept}
\begin{itemize}
\item \textbf{Dvija (twice-born):} Brahmins, Kshatriyas and Vaishyas receive the sacred thread --- they are born once physically and once spiritually.
\item \textbf{Ekaja (once-born):} the Shudra does not receive the thread --- born only once.
\end{itemize}
\end{KeyConcept}

\begin{itemize}
\item Under the Varna model, the three Dvija are the \textbf{Savarna} (noble Varna); the Shudra, and beyond the four a fifth, form the rest.
\item The \textbf{Pancham Varna (fifth Varna)} is the untouchable --- a Varna with no Varna.
\end{itemize}

\begin{QuickRevision}
\begin{itemize}
\item Ambedkar: the original kings were Shudras; denied the sacred thread, they became the lowest.
\item Sacred thread = Upanayana (spiritual rebirth).
\item Dvija (twice-born: B, K, V) vs Ekaja (once-born: Shudra).
\end{itemize}
\end{QuickRevision}

\section{Racial Theory of Caste}

The racial theory was given by the \textbf{British}, because the idea of race came from Europe. When the British decided to understand the population, they began classifying it (Mongoloid, Caucasoid, Negroid) --- race being the first category ``scientifically'' classified. They argued that \textbf{caste is nothing but race}.

\begin{DataFact}
Proponents of the racial theory: \textbf{Herbert Hope Risley} (director of the first ethnographic department of India; one of the first census commissioners), \textbf{D. N. Majumdar}, and \textbf{G. S. Ghurye}.
\end{DataFact}

\begin{CriticalPoint}
You cannot understand Ghurye without the racial theory --- it is the basis of his first book on caste.
\end{CriticalPoint}

The story: the \textbf{Aryans} (fair-skinned, pointed nose) came from outside and clashed with the \textbf{indigenous people} (dark-skinned, snub nose) of the Indus Valley (Harappan) civilisation. From the clash of these two races came \textbf{Vedic Hinduism}, and from Vedic Hinduism came caste.

\begin{itemize}
\item Vedic Hinduism did not exist before the Aryans; the written record of Sanskrit came with them. The Aryans are thus regarded as the makers of Hinduism.
\item The Aryans brought a \textbf{threefold division} already practised in Persia --- priestly class, warrior class and merchants --- and reproduced it in India as Brahmin, Kshatriya and Vaishya. The indigenous people became the Shudras.
\item The Indus Valley civilisation is generally regarded as an \textbf{egalitarian society} with a single system --- no caste structure.
\end{itemize}

\begin{KeyConcept}
To preserve racial purity, the Aryans practised \textbf{endogamy} (marriage within the same group) and \textbf{restricted hypergamy} (Anuloma). \textbf{Hypergamy} = an upper-caste man marrying a lower-caste woman; \textbf{hypogamy} (Pratiloma) = an upper-caste woman marrying a lower-caste man (prohibited).
\end{KeyConcept}

Endogamy and restricted hypergamy became instruments to preserve, conserve and maintain Aryan culture (language + culture = race). They made both groups (Aryan and non-Aryan) \textbf{closed groups}, and this closed group later became what we call \textbf{caste}.

\begin{DataFact}
\begin{itemize}
\item \textbf{Varna is a Sanskrit word meaning ``colour''} --- the four Varnas are literally a colour spectrum, with the darkest-skinned placed lowest. This is the original racial theory of caste.
\item Aryan invasion timeline: the Harappan civilisation ended around \textbf{1700 BC}; around \textbf{1500 BC} the Aryan nomads took control of the Indus Valley; \textbf{500 BC} marks the end of Aryan rule (and the start of Magadha rule).
\item The Aryans are seen as \textbf{cattle herders} --- nomadic warrior groups; \textbf{cattle were a sign of status and wealth} (today's equivalent of a big car). Later cattle became so important that it was made illegal to kill or eat them. The caste system officially ended in \textbf{1950 AD}.
\end{itemize}
\end{DataFact}

\begin{CriticalPoint}
The Aryan-invasion explanation for the Harappan decline (as found in R. S. Sharma's NCERT account) is regarded as \textbf{one of the weakest reasons} for the decline of the Harappan civilisation.
\end{CriticalPoint}

\begin{QuickRevision}
\begin{itemize}
\item Racial theory (Risley, Majumdar, Ghurye): Aryans (fair, pointed nose) vs indigenous (dark, snub nose).
\item Clash of races gives Vedic Hinduism, which gives caste.
\item Endogamy + restricted hypergamy make closed groups, which become caste.
\item Varna = ``colour''; Aryan timeline: 1700 BC (Harappan end), then 1500 BC (Aryan rule), then 500 BC (Magadha).
\end{itemize}
\end{QuickRevision}

\section{Occupational Theory of Caste}

This theory was given by \textbf{Nesfield} and \textbf{Denzil Ibbetson}. Nesfield accepts the racial theory but argues it is too broad to explain how caste actually formed; he therefore asks what lay beyond race.

\begin{KeyConcept}
Nesfield: \textbf{caste developed as a natural outcome of the division of labour.}
\end{KeyConcept}

\begin{itemize}
\item The people defeated by the Aryans were regarded as \textbf{Dasas (slaves)}; the Aryans, with their skills and discipline, became the masters.
\item With different occupations came different caste groups, and these built \textbf{guilds} around their occupation --- the guild system (given by Ibbetson; the guild is the \textbf{Shreni}).
\item Initially, caste was \textbf{not fixed or hereditary}; no occupation was superior or inferior --- every occupation was functional and necessary for the economy.
\item Gradually, the \textbf{Brahmins} (priests, teachers, advisors) gained respect: they were always seen with kings, and kings depended on them --- giving the impression that being Brahmin was being superior.
\end{itemize}

\begin{ExampleBox}
The dependency factor is key: the doctor is the most important person in a hospital because the whole staff depends on the doctor. Occupations with many dependents became high-status; those easily available became low-status.
\end{ExampleBox}

\begin{QuickRevision}
\begin{itemize}
\item Occupational theory (Nesfield, Ibbetson): caste = natural outcome of the division of labour.
\item Dasas (slaves); guilds (Shreni); hereditary occupation.
\end{itemize}
\end{QuickRevision}

\section{Political (Brahmanical) Theory}

The political theory is also called the \textbf{Brahmanical theory}. It holds that the caste system is \textbf{the making of the Brahmins} --- Brahmins manufactured the system to manipulate people and gain power.

\begin{QuickRevision}
\begin{itemize}
\item Political/Brahmanical theory: caste = the making of the Brahmins (to gain power).
\end{itemize}
\end{QuickRevision}

\section{The Two Most Important Theories}

\begin{CriticalPoint}
The two most important theories are the \textbf{Racial} and the \textbf{Occupational} theories, because their features turn out to be the features of caste itself:
\begin{itemize}
\item \textbf{Racial:} endogamy and restricted hypergamy.
\item \textbf{Occupational:} hereditary (occupation passed down).
\end{itemize}
\end{CriticalPoint}

\begin{QuickRevision}
\begin{itemize}
\item The two important theories: Racial (endogamy + hypergamy) and Occupational (hereditary).
\end{itemize}
\end{QuickRevision}

\section{Senart's Caution: Caste vs Varna}

\textbf{Senart} argues that all theories of origin confuse \textbf{caste with Varna} --- they look at caste through the Varna lens. The word \textbf{caste is alien} (Portuguese); India's indigenous word is \textbf{Jati}. Jati and Varna are indigenous; ``caste'' is foreign --- hence the persistent confusion.

\textbf{M. N. Srinivas} wrote the article \textbf{``Varna and Caste''} to make the distinction clear.

\begin{QuickRevision}
\begin{itemize}
\item Senart: don't confuse caste (alien, Portuguese) with Varna (indigenous); jati is indigenous.
\end{itemize}
\end{QuickRevision}

\section{Varna Model vs Jati Model (M. N. Srinivas)}

Srinivas argues that the fourfold classification (Brahmin, Kshatriya, Vaishya, Shudra) comes from the Sanskrit texts (Vedas, Manu Smriti) --- it is purely \textbf{book view}. The moment one steps outside the book into the field, one does not find these four; one finds actual \textbf{communities (jatis)}.

\begin{center}
\begin{tabularx}{\textwidth}{YYY}
\rowcolor{AccentDark}
\thead{Attribute} & \thead{Varna Model} & \thead{Jati Model} \\
Basis & Book view --- ancient Sanskrit texts (Vedas, Manu Smriti) & Field view --- actual communities observed in villages \\
Number of categories & 4 (Brahmin, Kshatriya, Vaishya, Shudra) & About 4000 castes all over India \\
Mobility & None --- fixed, immutable, hierarchically arranged & Mobility exists --- caste is mutable and nebulous \\
Hierarchy & Fixed and definite & ``Muddled'' --- regionally variable (Dipankar Gupta) \\
Scale & Macro-cosmic view; an ideal type / model & Local, empirical, region-specific \\
\end{tabularx}
\end{center}

Srinivas is regarded as the \textbf{first Indian sociologist} to understand India not through books but through the field --- using field work and \textbf{participant observation}. This was a paradigm shift in the sociology of India.

\subsection{The Rampura village study}

\begin{itemize}
\item Srinivas found not four but \textbf{many} categories, e.g.\ the \textbf{smith} caste.
\item Even among Brahmins there was hierarchy --- Hoysala Brahmin, Madhav Brahmin, Krishna Vaishnav Brahmin.
\item The distinction between \textbf{kachcha (raw) food} and \textbf{pakka (cooked) food}: refusing a person's cooked food marks them as lower; accepting it marks them as equal or higher.
\item The smith would accept cooked food only from the \textbf{Hoysala Brahmin} (so the smith saw himself as their equal), but the Hoysala Brahmin never reciprocated --- and other castes saw the smith as even more untouchable than the untouchables.
\end{itemize}

\begin{KeyConcept}
\textbf{Hierarchy is not fixed.} Even if Brahmin is superior and the untouchable inferior, the \textbf{middle region of the hierarchy is highly vague} --- Dipankar Gupta calls it a \textbf{``muddled hierarchy.''}
\end{KeyConcept}

\subsection{Mobility in caste}

Srinivas found mobility in caste. Caste is therefore \textbf{mutable and nebulous}, not immutable. The key is the \textbf{scale} on which we measure:

\begin{ComparisonBox}
\begin{itemize}
\item \textbf{Ritual scale} (purity--impurity): mobility does not occur --- a Shudra remains ritually impure even after becoming an IAS officer.
\item \textbf{Secular/merit scale} (individual achievement): mobility does occur --- the Shudra who becomes an officer has risen, even if not ritually.
\end{itemize}
\end{ComparisonBox}

Mobility in caste comes from political and material conditions. For example, in U.P. the \textbf{Ahir} are considered dominant because they have political and material power.

\begin{CriticalPoint}
On ``is caste a myth and class a reality?'' --- the answer given is that \textbf{class is also a myth}: class means an open system, but nothing is absolutely open. Models (Varna, class) are what we read; reality is far more complex and can never be academically presented in neat, chapter-wise form.
\end{CriticalPoint}

\begin{ExampleBox}
In the field one caste may hold \textbf{multiple occupations} and \textbf{multiple castes} may share one occupation (as F. G. Bailey found at Bisipara, Odisha) --- so the occupational link is not fixed.
\end{ExampleBox}

\begin{QuickRevision}
\begin{itemize}
\item Varna model (book, 4 castes, no mobility) vs Jati model (field, ~4000 castes, mobile, ``muddled'').
\item Srinivas = first field-view / participant-observation sociologist.
\item Rampura study; kachcha vs pakka food; ritual vs secular scale of mobility.
\end{itemize}
\end{QuickRevision}

\section{Varna Vyavastha \& Related Concepts}

\textbf{Varna Vyavastha} is the term scholars use for the Varna system. Beyond the four Varnas there is the \textbf{Pancham Varna} (fifth) --- the untouchables.

\begin{KeyConcept}
\begin{itemize}
\item \textbf{Savarna:} the four ``noble'' Varnas (B, K, V, S).
\item \textbf{Avarna:} those outside --- the untouchables (no Varna; the ``impure'' Varna).
\item \textbf{Dvija:} twice-born (the three that receive the sacred thread).
\item \textbf{Ekaja:} once-born (Shudra and Avarna). Shudra and Avarna are both Ekaja, which brings them very close to each other.
\end{itemize}
\end{KeyConcept}

Even within Brahmins there is hierarchy --- \textbf{Trivedi} (read three Vedas), \textbf{Chaturvedi} (four Vedas), \textbf{Dvivedi} (two Vedas). There are even \textbf{untouchable Brahmins}: the priests who perform \textbf{death rituals} are seen as polluted by the ``pure'' priests.

\begin{ExampleBox}
Occasions of purity (birth/Naamkaran, havan, new house, wedding) bring a Brahmin priest in; but the priest who handles death is treated as polluted. There is even specialisation among Pandits --- some exorcise spirits, some handle death, some handle weddings.
\end{ExampleBox}

\subsection{Who are the Dalits?}

\begin{KeyConcept}
The \textbf{Pancham Varna (untouchables)} includes three book-view categories: \textbf{Dalits}, \textbf{Chandal}, and \textbf{Malechha}.
\end{KeyConcept}

\begin{itemize}
\item \textbf{Chandal} --- the child born of a prohibited union (an upper-caste woman and a lower-caste man); such a child has \textbf{no caste} --- a punishment designed to prevent intermixing and preserve caste purity.
\item \textbf{Malechha} --- outsiders/foreigners, seen as impure.
\item Jati is a community term --- we speak of Muslim jaat, Hindu jaat, Sikh jaat, Bengali jaat, Maratha jaat, Yadav jaat, Kunbi jaat, Oraon jaat, Munda jaat (tribals too are communities/jatis).
\end{itemize}

\begin{QuickRevision}
\begin{itemize}
\item Varna vyavastha; Pancham Varna (untouchables).
\item Savarna vs Avarna; Dvija vs Ekaja; untouchable Brahmins (death-ritual priests).
\item Pancham Varna = Dalits, Chandal, Malechha.
\end{itemize}
\end{QuickRevision}

\section{Dalits, Scheduled Castes \& OBCs}

\subsection{Scheduled Castes}

\textbf{Scheduled Caste (SC)} is a \textbf{government-notified} category (listed in the Gazette of India). \textbf{Not all Dalits are SCs} --- only those on the list are.

\begin{ExampleBox}
Like the Below Poverty Line list, the SC list has \textbf{inclusion and exclusion errors} --- governments, being human, put the wrong people in and leave the right people out (the analogy of the friend who ruins the party).
\end{ExampleBox}

\begin{DataFact}
Dalits are \textbf{not exclusive to Hindus} --- there are Dalit Christians, Buddhists, Jains, Sikhs and Muslims. But the SC category is given only to \textbf{Hindu, Sikh and Buddhist Dalits} --- \textbf{not to Muslim or Christian Dalits} (Presidential Order, 1950). Sikhism and Buddhism are treated as Hindu religions, so their Dalits qualify.
\end{DataFact}

A court case is ongoing over SC status for Muslim/Christian Dalits.

\subsection{Dalit elites}

Among Dalits there are \textbf{Dalit elites}. Because SC status brings reservation, an SC holder (say, a Hindu Dalit) rises while his Muslim-Dalit friend does not. The benefit makes him an officer, builds \textbf{social capital}, and can carry him to an MLA candidacy.

\begin{KeyConcept}
It is the \textbf{intersectionality of religion and caste} that produces mobility among individuals. Sociological analysis focuses on \textbf{structures}, not individuals --- like a \textbf{geologist} who explains earthquakes through tectonic plates, not individual whim.
\end{KeyConcept}

\subsection{OBCs and the ``backward-forward''}

Shudra is also not homogeneous: there is the \textbf{forward Shudra (upper caste)} and the \textbf{backward Shudra (lower caste)}. The \textbf{OBC (Other Backward Class)} category is for the backward Shudras.

\begin{GovtPolicy}
The Constitution speaks of ``other socially and educationally backward classes''; the \textbf{Mandal Commission (1979)} was set up to identify them and drew up criteria and a list. ``Other'' means other than the Scheduled Castes and Scheduled Tribes. Within the OBC category, the better-off form the \textbf{creamy layer}, who are meant to be excluded from the benefits.
\end{GovtPolicy}

\begin{ExampleBox}
The same Jat is OBC in one state and forward in another --- \textbf{regional variation} makes the book view (``Jat'') useless; it is a field-view category.
\end{ExampleBox}

\begin{CriticalPoint}
\begin{itemize}
\item A new \textbf{Varna Vyavastha} has been created by the government: \textbf{General, EWS, OBC, SC, ST} --- and none of these is homogeneous.
\item \textbf{Satish Deshpande} coined \textbf{``backward-forward''} --- backward as Shudra, forward within the Shudra. He says the \textbf{Jat quota demand is more an expression of caste entitlement than a claim to class disadvantage.}
\end{itemize}
\end{CriticalPoint}

\begin{QuickRevision}
\begin{itemize}
\item SC = government-notified list; not all Dalits are SCs.
\item SC only for Hindu/Sikh/Buddhist Dalits (Presidential Order 1950); not Muslim/Christian.
\item Dalit elites; OBC (Mandal 1979) with creamy layer; backward-forward (Deshpande).
\end{itemize}
\end{QuickRevision}

\section{Caste Census}

India never had ``caste'' --- it had \textbf{Jati and Varna}. The British did not know these; they only asked ``who are we?'' and, to answer, conducted a \textbf{census} (a survey).

\begin{DataFact}
The first census was conducted in \textbf{1871} (initially the North-Western Province, becoming an all-India census around 1891--1901).
\end{DataFact}

A census is one form of survey. The British wanted a \textbf{``neat and clean''} picture --- tables, rows and columns --- so they \textbf{quantified a qualitative reality}, dividing people into caste, religion and language.

\begin{itemize}
\item They referred both to \textbf{ancient Brahmanical texts} and to \textbf{vernacular texts} (local/folk/tribal literature, jati puranas) --- because tribal life is not in the Vedas.
\item The result was a bewildering list of categories --- Kunbi, Yadav, Baidya, Bhargava, Jat, Suthar, Gujjar, Muslim, Buddhist, Jain, and so on.
\item The problem became ``\textbf{who amongst them is caste?}'' --- people did not know their own caste; in Punjab the census commissioners found people could not say. Caste was \textbf{imposed on India by the British} and thereby manufactured.
\end{itemize}

\begin{CriticalPoint}
The caste census did not classify the population --- it \textbf{divided} us into castes.
\end{CriticalPoint}

\subsection{Nicholas Dirks and Foucault}

\begin{KeyConcept}
\begin{itemize}
\item \textbf{Nicholas Dirks} (``Castes of Mind''): caste is the \textbf{product of the colonial project of classifying people}; the census was a \textbf{``cultural technology of power.''}
\item ``Colonial conquest not only creates knowledge, but is produced by it as well.''
\item \textbf{Michel Foucault:} knowledge is power --- knowledge cannot be separated from power.
\end{itemize}
\end{KeyConcept}

\begin{DataFact}
The last SC/ST caste census was in \textbf{1931}. The recent debate was over an \textbf{OBC census} --- the enumeration of Other Backward Classes, not of Dalits or tribals.
\end{DataFact}

\begin{QuickRevision}
\begin{itemize}
\item First census 1871; caste imposed by the British (manufactured).
\item Nicholas Dirks: census = ``cultural technology of power''; Foucault: knowledge is power.
\item Last SC/ST caste census: 1931; recent debate = OBC census.
\end{itemize}
\end{QuickRevision}

\section{Orientalism, Indology \& Occidentalism}

\subsection{Orientalism}

\begin{KeyConcept}
\textbf{Orientalism} = the representation of non-Western people by Western people.
\end{KeyConcept}

To the Orientalist, India is caste, snake charmers and priests --- an image the film \textbf{Slumdog Millionaire} reinforced even among people who had never known India.

\begin{CriticalPoint}
But not all Europeans represented India badly. \textbf{Voltaire} loved the Ganga; \textbf{William Jones} found similarities between Sanskrit and English; \textbf{Max M\"uller} coined ``Aryan brotherhood.'' \textbf{Trotman} calls these lovers of India \textbf{Indomaniacs}, and the denigrators \textbf{Indophobics}.
\end{CriticalPoint}

With the consolidation of British rule, the mood turned from \textbf{Indomania to Indophobia} --- the \textbf{European U-turn}.

\begin{ComparisonBox}
\begin{itemize}
\item \textbf{William Jones} is an \textbf{Indologist} because he established the first department dedicated to understanding Indian culture through the ancient Brahmanical texts --- the founder of Indology.
\item \textbf{Travellers} (Albiruni --- ``Kitab-ul-Hind''; Abdul Razak Samarkandi --- Vijayanagar; Fran\c{c}ois Bernier and Tavernier --- the Mughal empire) studied \textbf{empires}, not Indian culture.
\end{itemize}
\end{ComparisonBox}

\subsection{Occidentalism}

\begin{KeyConcept}
\textbf{Occidentalism} is the exact opposite of Orientalism --- the non-West representing the West (e.g.\ Nehru's depiction of the West; the ``Great Indian wedding''; Monsoon Wedding; Bride and Prejudice; the Bollywoodisation of South Indian cinema and the Hollywoodisation of Bollywood --- a pattern of homogenisation).
\end{KeyConcept}

\begin{QuickRevision}
\begin{itemize}
\item Orientalism = West representing non-West; Occidentalism = non-West representing West.
\item Indomania to Indophobia = the European U-turn.
\end{itemize}
\end{QuickRevision}
